\section{Introduction}
\label{sec:introduction}

Variational formulations turn a broad class of elliptic boundary-value problems into convex energy
minimization. Classical finite elements approximate the admissible space by locally supported
polynomials, and the relation between polynomial degree, mesh scale, and degrees of freedom is well
understood \citep{BrennerScott2008}. Neural and ridge-function methods instead construct global
nonlinear trial spaces, so that the approximation problem concerns the selection and
reoptimization of atoms rather than the design of a mesh. The deep Ritz method illustrates the
reach of variational neural solvers \citep{EYu2018}. The present work analyzes a deterministic
greedy selection rule for such nonlinear elliptic energies.

Orthogonal greedy approximation (OGA) has a sharp Hilbert-space theory. Under suitable source and
entropy conditions it produces optimal approximation rates
\citep{SiegelXuOGA2022,LiSiegelEntropy2024}. Chebyshev-type greedy methods extend the basic descent
argument to convex energies in Banach spaces \citep{TemlyakovCGA2013,TemlyakovGreedy2014}, and
biorthogonal or inexact variants show how selection and correction errors enter the energy recursion
\citep{DereventsovTemlyakov2022}. These results supply the approximation theory on which the analysis
below relies. Applying them to a nonlinear partial differential equation (PDE) whose atoms are drawn
from a finite candidate pool also requires accounting for numerical quadrature, finite dictionary
approximation, and inexact coefficient minimization. The three effects are distinct: quadrature
perturbs the evaluated score, a finite pool replaces the continuous supremum by a maximum over the
pool, and an active-space solver returns only an approximate minimizer.

The dictionary used below consists of shallow \(\relu^k\) ridge functions. Here
\(\relu^k(s)=[s]_+^k\), with \([s]_+=\max\{s,0\}\), and the domain is a bounded subset of
\(\R^d\). Greedy neural training has been used successfully for PDEs
\citep{SiegelHongJinHaoXu2023}. A pool that is drawn once and then held fixed during the greedy
solve provides a way to formulate weak selection, by which we mean selection of an atom whose score is a
prescribed fraction of the best score available in the full dictionary:
\cite[Theorem~3.1]{XuXu2025} prove a convex-energy bound under Hilbert smoothness and
source/entropy assumptions, while their Theorems~5.1 and 5.2 quantify dictionary sizes that
realize weak selection. Their analysis assumes exact coefficient minimization over the active
span. The finite-pool inequality below separates the correction error from the score and coverage
errors. The universality of such constructions for nonlinear convex variational problems is
examined by \cite{BernaFalco2025}. Variation spaces and sharp entropy estimates for shallow
networks provide the source and width information used in the conditional estimates
\citep{SiegelXuVariation2023,SiegelXuSharp2024,MaoSiegelXu2026}.

The analysis has three parts. First, a one-step descent estimate separates pool
coverage, score quadrature, radial stationarity, and coefficient-correction errors.
Second, projection bounds for the computed nested spaces retain the source remainder
and the observed orthogonal increments; an exact quartic expansion connects these
bounds to nonlinear energy errors. Third, Section~\ref{sec:c5-hessian} gives a local
Hessian comparison for a concrete semilinear family and identifies the relative-score,
invariant-neighbourhood, and entropy assumptions needed for a conditional convergence
rate. The entropy estimate concerns the same projected, unrenormalized restarted
dictionary throughout. The numerical results evaluate the finite-trajectory bounds
and selected transfer quantities; they do not establish an unconditional asymptotic rate.

The convex Chebyshev algorithms of \cite{TemlyakovCGA2013,TemlyakovGreedy2014} and the inexact
weak-greedy framework of \cite{DereventsovTemlyakov2022} are closest to the present setting. The
former supplies the generic descent argument, and the latter shows how approximate selection and
correction enter an energy recurrence. The Hilbert OGA results of
\cite{SiegelXuOGA2022,LiSiegelEntropy2024} apply after a local transfer to a fixed Hessian metric;
they do not by themselves control a nonlinear PDE score or selection from a finite pool.
\cite{XuXu2025} quantify deterministic and randomized finite dictionaries under weak selection, but
they use exact active-space minimization and do not analyze a degenerate \(p\)-energy. The present
work examines the additional conditions needed for nonlinear scores, finite pools, and inexact
active-space minimization.

For \(p\)-growth energies the metric that matches the energy is not the ambient Sobolev norm. The
nonlinear map
\[
  V(\xi)=|\xi|^{(p-2)/2}\xi
\]
defines the natural distance, as in the finite-element analysis of the \(p\)-Laplacian
\citep{BarrettLiu1993,DieningKreuzer2008,LeePark2025}. A squared \(V\)-distance has the scale of an
energy gap, so an energy bound of order \(m^{-\alpha}\) gives a \(V\)-distance bound of order
\(m^{-\alpha/2}\) and, for \(p\geq2\), a \(W^{1,p}\)-error bound of order \(m^{-\alpha/p}\). This
conversion has the same form for every \(p\), although the energy exponent may depend on \(p\). We
prove the equivalence for the standard pure, regularized, and reaction \(p\)-energies and apply the
corresponding conversion to each of them.

The atomic source condition of \Cref{ass:atomic} and the directional smoothness estimate of
\Cref{ass:directional} apply to convex energies and yield an \(O(m^{-1})\) energy bound for the ideal
continuous CGA. The same arguments produce an explicit one-step inequality for finite-pool CGA, in
which four contributions appear separately: continuous-dictionary coverage, score quadrature, radial
active-space stationarity, and coefficient optimization.

Two conditional refinements use additional residual or Hessian information. The residual-comparison
estimate combines power-type convexity with a lower bound for the dictionary dual seminorm in terms
of the full dual norm. A local Hessian transfer uses coercivity, a dictionary norm comparison, and a
relative score condition to obtain a weak-OGA estimate in a frozen Hilbert metric. Neither set of
assumptions implies the other, and both are stated as conditional estimates.

The finite trajectory produced by the algorithm can also be analyzed without requiring uniform
Hessian coercivity. For a fixed finite symmetric dictionary and a fixed starting index \(m_0\), an
explicit approximation of the source with coefficient budget \(B\) and remainder \(\sigma\) gives a
projection estimate
\[
 r_{m_0+j}^2\le \sigma^2+\frac{4B^2}{\mathcal W_j},
 \qquad
 \mathcal W_j=\sum_{m=m_0}^{m_0+j-1}\frac{\underline\theta_m^2}{\ell_m^2}.
\]
Here \(r_m\) is the norm of the projection error at state \(m\), \(\ell_m\) is the innovation norm of
the selected candidate, and \(0\leq\underline\theta_m\leq\theta_m\) records the selected direction
against the best score available in the frozen dictionary. For the one-dimensional \(p=4\) energies,
whose density is quartic in the gradient, an exact expansion connects this projection error to the
energy gap. The estimate applies to the active spaces actually generated and carries a source
remainder; it does not assume that the target lies in the finite dictionary. The stronger
initial-value form of the bound is stated and evaluated numerically. A second estimate uses the
normalized innovation \(\gamma_m=\theta_m S_m/\ell_m\). If
\(\gamma_m^2\geq c_0r_m^{2+2/\alpha}\), inversion of the exact projection recurrence bounds
\(r_{m_0+j}^2\) by
\(\big(r_{m_0}^{-2/\alpha}+c_0j/\alpha\big)^{-\alpha}\). The same quartic comparison
then gives the corresponding energy and metric bounds. Each condition is evaluated from the retained
projection drops; on a finite index range, positivity alone does not single out a unique exponent.

The numerical study evaluates the estimates together with the conditions on which they rest. Eight
representative cases cover linear reaction--diffusion, cubic and hyperbolic-sine semilinear
energies, and pure, regularized, and reaction \(p\)-energies. Five of them compare finite-pool CGA, a
random-feature model (RFM), and finite elements (FEM) at equal numbers of independent coefficients.
Random features are used in their fixed-feature sense \citep{RahimiRecht2007}; conditioning of
linearized neural trial spaces is discussed by \cite{ChenChiEYang2022} and
\cite{MaoXuXu2026}. The comparison concerns coefficient counts and uses a common evaluator. It does
not attribute an unproved convergence theorem to the compared methods and does not compare
wall-clock efficiency. P1 and P3 elements provide low- and high-order mesh-based references, and P2
results are retained in the supplementary material.

For the finite-trajectory estimates we use the one-dimensional pure and regularized $p=4$
trajectories. The calibration range is $N=8$--$32$, and the subsequent range uses the same stored
prefixes from $N=33$ to $64$ with the calibration constants fixed in advance. The reported
fractions record the number of steps satisfying the stated sufficient condition on each finite
range. They provide numerical evidence on the displayed trajectory; they do not verify continuous-
dictionary coverage or an entropy estimate uniformly over increasing ranges. The approximate
source is constructed once at the starting index from the fixed candidate and auxiliary reference
dictionary. The source remainder is small relative to the starting residual but is not negligible
at the terminal scale; it therefore limits the resolution of the source upper bound without giving
a lower bound on the actual algorithmic error. Neither calculation verifies the continuous-
dictionary coverage and entropy assumptions of the Hessian-transfer result, so the measured orders
and fitted innovation powers are reported as finite-range observations.

The remainder is organized as follows. \Cref{sec:setting} defines the variational setting,
dictionaries, and algorithms. \Cref{sec:c5-descent} proves the basic descent estimates, and
\Cref{sec:pgeometry} develops the energy families and natural-distance conversion for \(p\geq2\).
\Cref{sec:c5} combines residual comparison, Hessian transfer, finite-trajectory
bounds, and fractional innovation. \Cref{sec:numerics} tests these estimates on the computed
trajectories; \Cref{sec:results} presents equal-DOF comparisons, and \Cref{sec:discussion}
summarizes the conclusions and their scope.
